\section{Worked Example: Canonical Wamura Scenario}
\label{sec:worked_example}

We instantiate the protocol on a canonical Wamura scenario
adapted from~\citep[Discussion]{lasser2026horizon}, showing the
behaviour under both a false risk claim and a genuine risk.

\subsection{Setup}

A newly discovered capability tuple $(d_1, d_2)$ is flagged by
agent $a_j$ as having a potential exercise pathway.  Agent $a_j$
assigns $\pR(0) = 0.7$ (moderately high confidence that the
pathway is dangerous).  The preemptive restriction
criterion~\citep[Proposition~1]{lasser2026horizon} triggers
because $\pR(0)$ exceeds the threshold, and keystone $d_1$ is
restricted.  We analyse both ground truths:
\begin{itemize}
\item $H_0$ (pathway is harmless; $\theta = 0$),
\item $H_1$ (pathway is genuinely dangerous; $\thetaone = 1$).
\end{itemize}

\paragraph{Parameters.}
\begin{itemize}
\item Prior: $\pR(0) = 0.7$ (probability of $H_1$).
\item Alternative hypothesis: $\thetaone = 1$ (``if the pathway
  is dangerous, it activates every time the capability is
  exercised'').
\item Aggregate per-step detection probability: $\deltaagg = 0.95$
  (the multi-channel architecture detects activation with $95\%$
  per-step probability).
\item Test duration: $\testdur = 3$ steps.  The
  $\testopt$ bound (Proposition~\ref{prop:tau_opt}) gives
  $\testopt = \lceil -1/\log(1 - 0.95) \rceil = \lceil 0.334
  \rceil = 1$; we use $\testdur = 3$ for additional sensitivity
  margin.  Running $\testdur > \testopt$ reduces $\beta$
  further (each extra step multiplies the miss probability by
  $1 - \deltaagg$) at the cost of extending the damage-exposure
  window, a conservative trade-off that the worked example
  illustrates but does not optimize.
\item Window false-negative rate:
  $\beta(3) = (1 - 0.95)^3 = 1.25 \times 10^{-4}$.
\item False-positive rate: $\alphafp = 0.01$.
\item Damage tolerance: $\damagebound = 0.02 \cdot \volL$ ($2\%$
  of total capability volume).
\end{itemize}

\subsection{Without controlled relaxation (the Wamura problem)}

Under the standard framework without controlled relaxation,
Proposition~\ref{prop:asym_evidence}(b) gives $\pR(t) = 0.7$ for
all $t$.  The restriction $\Xirestrict$ remains in force
indefinitely; $\volL$ is permanently reduced by the restriction
cost $|\Delta \volL(\text{restrict } d_1)|$ for the rest of time.
The ground truth never surfaces; the $70\%$ estimate persists
forever.

\subsection{With controlled relaxation, ground truth $H_0$}

\paragraph{Test~1.}
Initiation criteria met (restriction is stale,
$\IV > \damagebound \cdot \pR$).  Scope $\scope$ constructed
around $\{d_1, d_2\}$ with $\maxcon(\scope) = 0.015 \cdot \volL <
\damagebound$.

Outcome: no contraction observed.  Under $H_0$, the pathway
never activates, so $O_n = 1$ only via a false positive:
$\Pr(O_n = 1 \mid H_0) = \alphafp = 0.01$; hence
$\Pr(O_n = 0 \mid H_0) = 0.99$.  Using
Equation~\eqref{eq:lik_h1} with $\thetaone = 1$:
$$
\Pr(O_n = 0 \mid H_1)
   = 1 \cdot \beta + 0 \cdot (1 - \alphafp)
   = 1.25 \times 10^{-4}.
$$
The likelihood ratio is
$$
\Lzero = \frac{1.25 \times 10^{-4}}{0.99}
       \approx 1.263 \times 10^{-4}.
$$
Applying~\eqref{eq:update_neg},
$$
\pR(t_1) = \frac{0.7 \cdot 1.263 \times 10^{-4}}
                {0.7 \cdot 1.263 \times 10^{-4} + 0.3}
        \approx 2.95 \times 10^{-4}.
$$
A single test with no contraction drops $\pR$ from $0.7$ to
$\approx 3 \times 10^{-4}$.  The dramatic drop reflects the high
detection sensitivity: under $H_1$ with $\thetaone = 1$, the
detector would have caught the activation with probability
$1 - \beta = 99.99\%$.  Absence of detection is extremely
unlikely under $H_1$ and expected under $H_0$.

\paragraph{Test~2 (confirmation).}
$\pR(t_1) \approx 3 \times 10^{-4}$ is already below the typical
restriction-firing threshold $p_{\mathrm{thresh}}$ (the
derived-quantity scalar at which the ex-ante cost-benefit
inequality of~\citep[Proposition~1]{lasser2026horizon} first
fires; see Appendix~\ref{app:proofs},
Proposition~\ref{prop:monotone_overcaution} for the derivation),
so the restriction is lifted.  If the meta-policy
calls for a confirmation test,
$$
\pR(t_2) = \frac{\pR(t_1) \cdot \Lzero}
                {\pR(t_1) \cdot \Lzero + (1 - \pR(t_1))}
        \approx 3.7 \times 10^{-8}.
$$

Equation~\eqref{eq:n_clear} with the worked-example parameters
($\Lzero \approx 1.263 \times 10^{-4}$, $\pR(0) = 0.7$,
$p_{\mathrm{clear}} = 0.01$) gives
$$
N_{\mathrm{clear}}
  = \left\lceil
    \frac{\log\bigl(0.01 \cdot 0.3 /(0.7 \cdot 0.99)\bigr)}
         {\log 1.263 \times 10^{-4}}
  \right\rceil
  = \lceil 0.606 \rceil = 1 \text{ test}.
$$
A single test suffices because the per-test likelihood ratio is
so extreme (the detector would catch a real pathway with
$99.99\%$ probability).  The worked example confirms this:
after one test, $\pR$ is already $3 \times 10^{-4}$.  The
second test is a confirmation that drives $\pR$ to negligible
levels.

\subsection{With controlled relaxation, ground truth $H_1$}

Suppose instead the ground truth is $H_1$: the pathway is
genuinely dangerous with $\thetaone = 1$.  For test~1 with
$\testdur = 3$,
$$
\Pr(O_n = 1 \mid H_1)
   = 1 \cdot (1 - 1.25 \times 10^{-4})
     + 0 \cdot 0.01
   = 0.999875.
$$
The test detects contraction with near-certainty.  The
likelihood ratio is
$$
\Lone = \frac{0.999875}{0.01} = 99.9875.
$$
Applying~\eqref{eq:update_pos},
$$
\pR(t_1) = \frac{0.7 \cdot 99.9875}
                {0.7 \cdot 99.9875 + 0.3}
        = \frac{69.991}{70.291}
        \approx 0.9957.
$$
The risk estimate increases from $0.7$ to $0.996$, and the
restriction is reinforced.  Maximum damage during the test was
bounded by $\damagebound = 0.02 \cdot \volL$, which the
cooperative-closure scope construction guarantees.

\subsection{Summary}

\begin{table}[h]
\centering
\small
\begin{tabular}{@{}p{3.6cm}p{2.0cm}p{2.2cm}p{2.2cm}p{3.0cm}@{}}
\toprule
\textbf{Scenario} & \textbf{Ground truth} & \textbf{Tests to resolve} & \textbf{Final $\pR$} & \textbf{Damage incurred} \\
\midrule
False alarm, no CR
  & $H_0$
  & $\infty$ (never resolves)
  & $0.70$
  & $0$ during test, but permanent $\volL$ loss from restriction \\
False alarm, with CR
  & $H_0$
  & $1$ ($+1$ confirmation)
  & $3.7 \times 10^{-8}$
  & $0$ (pathway did not activate) \\
True risk, with CR
  & $H_1$
  & $1$
  & $0.996$
  & $\leq \damagebound = 0.02 \cdot \volL$ (bounded by
    Theorem~\ref{thm:damage_bound}) \\
\bottomrule
\end{tabular}
\caption{Behaviour of controlled relaxation on the canonical
  Wamura scenario.  CR~=~controlled relaxation.}
\label{tab:worked_example}
\end{table}

Table~\ref{tab:worked_example} collects the three cases.  The
framework without controlled relaxation never resolves a false
alarm; with controlled relaxation, a single test (backed by a
$95\%$-per-step multi-channel detector) suffices to clear the
false alarm to $3 \times 10^{-4}$ posterior, and a single
detection suffices to confirm a genuine risk to $99.6\%$
posterior.  Both outcomes carry bounded damage.
